\section{Syntax}\label{sec:syntax}

Baltmeri syntax was decided by the same votes as its morphology, and it is deliberately plain.

\paragraph{SVO, no inversion.} All nine languages are SVO at rest \citep{dryer2013}; only Germanic has verb-second \citep{holmberg2015}, and it loses 3--1. There is no inversion, ever: a fronted time adverb leaves subject and verb in place, \emph{Home me plavisme} ``tomorrow we will swim'', where Swedish or German would invert.

\paragraph{Modifiers precede.} Adjective before noun, genitive before noun (\emph{Baltmerin kala} ``the sea's fish''), numeral before noun. Quantifiers take the plain plural (\emph{daug kalai} ``many fish''); the Finnic partitive lost.

\paragraph{Place before time.} Adverbial phrases go to the end of the clause, place before time (the Slavic and Germanic order; Finnic is flexible): \emph{odpočiv kaljuil suvil} ``rest on the rocks in summer''.

\paragraph{Existence.} Existence is expressed with \emph{er} and the adessive: \emph{Meril er jää} ``there is ice on the sea''.

\paragraph{Negation and questions.} \emph{Ne} precedes the verb, or the predicate after the copula; the yes/no particle \emph{li} follows the verb: \emph{Tu kohit li meria?} ``do you love the sea?''.

\subsection{Example sentences}
The Manual's three founding sentences, glossed. Abbreviations: \textsc{gen} genitive, \textsc{acc} accusative, \textsc{ine} inessive, \textsc{ade} adessive, \textsc{ela} elative, \textsc{all} allative, \textsc{pl} plural, \textsc{sup} superlative, \textsc{fut} future, \textsc{q} question particle.

\begin{exe}
\ex \gll Baltmeri-n kala er daug sort-i i dužus-i. \\
Baltic-\textsc{gen} fish is many sort-\textsc{pl} and size-\textsc{pl} \\
\glt `The fish in the Baltic Sea come in many varieties and sizes.'
\ex \gll Peleek-hülje-i odpočiv kalju-i-l suvi-l. \\
grey-seal-\textsc{pl} rest.3 rock-\textsc{pl}-\textsc{ade} summer-\textsc{ade} \\
\glt `Grey seals rest on the rocks in summer.'
\ex \gll Räim er pospolit-ist kala Baltmeri-se. \\
Baltic.herring is common-\textsc{sup} fish Baltic-\textsc{ine} \\
\glt `Herring is the most common fish in the Baltic Sea.'
\ex \gll Ja vid-u kala-a vesi-se. \\
I see-1\textsc{sg} fish-\textsc{acc} water-\textsc{ine} \\
\glt `I see a fish in the water.'
\ex \gll Tu koh-it li meri-a? \\
you love-2\textsc{sg} \textsc{q} sea-\textsc{acc} \\
\glt `Do you love the sea?'
\ex \gll Home me žegl-is-me Gotland-ile boot-se. \\
tomorrow we sail-\textsc{fut}-1\textsc{pl} Gotland-\textsc{all} boat-\textsc{ine} \\
\glt `Tomorrow we sail to Gotland in a boat.'
\ex \gll Ne plav-it talv-il! \\
not swim-2\textsc{sg} winter-\textsc{ade} \\
\glt `Don't swim in winter!'
\ex \gll Üks duž kala er duž-im kui tri mal-i. \\
one big fish is big-\textsc{comp} than three small-\textsc{pl} \\
\glt `One big fish is bigger than three small ones.'
\end{exe}

All eight sentences are reproduced verbatim by the engine from a hand-written analysis (\S\ref{sec:eval-sent}), with one exception discussed there.

\subsection{A text}
The Manual closes its grammar with a short text, given here in Baltmeri with a translation; new words it introduces are listed in Appendix~\ref{app:lexicon}.

\begin{quote}\small
\textbf{Baltmeri}\\
Baltmeri er ne glembok meri. Baltmerin vesi er ne müket solne, i talvil meril er jää. Daug kalai plav Baltmerise: räim, sild, tursk i lohi. Peleekhüljei ži saaril i kaljuil, i suvil de odpočiv päivel.

Kalajai lov räima sügisl. Vakar de vidil daug hüljei meril. Home me žeglisme Gotlandile, i ja vidsu dzintara rannal.

Ja kohu Baltmeria. Meri er zimne, ne glembok, i ne müket solne---ale meri er men.
\end{quote}
\begin{quote}\small\itshape
The Baltic is not a deep sea. The Baltic's water is not very salty, and in winter there is ice on the sea. Many fish swim in the Baltic: Baltic herring, herring, cod and salmon. Grey seals live on islands and rocks, and in summer they rest in the sun.

Fishers catch herring in autumn. Yesterday they saw many seals on the sea. Tomorrow we sail to Gotland, and I will see amber on the shore.

I love the Baltic. The sea is cold, not deep, and not very salty---but the sea is ours.
\end{quote}
